
\documentclass[11pt,a4paper]{article}
\usepackage[a4paper,margin=2.05cm]{geometry}
\usepackage[spanish,es-noquoting]{babel}
\usepackage[T1]{fontenc}
\usepackage[utf8]{inputenc}
\usepackage{lmodern,microtype}
\usepackage{amsmath,amssymb,amsthm,mathtools}
\usepackage{booktabs,longtable,array,enumitem}
\usepackage{xcolor,hyperref}
\pagecolor{white}\color{black}
\setlength{\parindent}{0pt}
\setlength{\parskip}{0.65em}
\newtheorem{teorema}{Teorema}
\newtheorem{observacion}[teorema]{Observación}
\newtheorem{conjetura}[teorema]{Conjetura}
\title{\bfseries N42 --- Coeficientes CKM como cruce con \(\pi\)\\
\large Recuperación estructural de \(28,-25/2,-2/3\)}
\author{HMT / auditoría técnica}
\date{}
\begin{document}
\maketitle

\begin{abstract}
N42 corrige el estatuto de N41: los coeficientes \(28,-25/2,-2/3\) no deben presentarse primero como una incógnita del species-pack. Aquí se reconstruyen desde la geometría de cruce \(\pi\)-Witt: la \(4\)-cara \(B_{\mathfrak C}=H_e\cap P_\pi\), el pivote \(\varphi\) \(p_\varphi=7\), el umbral de Steiner \(5\), la transición hexada-octada \(6\to8\) y el denominador trítico \(3\). El resultado reproduce la fórmula CKM de N41.
\end{abstract}

\section{Datos de cruce}

El núcleo de cruce de \(\alpha\) contiene:
\[
B_{\mathfrak C}=H_e\cap P_\pi,\qquad |B_{\mathfrak C}|=4.
\]
El pivote negativo de \(\varphi\) es:
\[
p_\varphi=7.
\]
El cierre de Steiner usa cinco puntos:
\[
p=5.
\]
La transición visible/ambiente es:
\[
W_{12}\to W_{24},\qquad 6\to8.
\]
El denominador de orientación es TRIT:
\[
3.
\]
La rama débil complementaria del cierre pentagonal se representa por:
\[
4.
\]

\section{Reglas de coeficientes}

\subsection*{Canal \(12\)}
El canal \(12\) usa dos copias de \(A\), resta el umbral de Steiner en \(h\), y añade la torsión del cruce \(B_{\mathfrak C}\times p_\varphi\):
\[
\boxed{
\theta_{12}=2A-5h+(4\cdot7)\Delta^\circ
=
2A-5h+28\Delta^\circ.
}
\]
El número \(28\) no es suelto: es el área de cruce \(4\times7\).

\subsection*{Canal \(23\)}
El canal \(23\) no lleva \(A\) directo; lleva el pivote \(\varphi\) en \(h\), y la torsión es el semicuadrado two-way del cierre pentagonal:
\[
\boxed{
\theta_{23}=7h-\frac{5^2}{2}\Delta^\circ
=
7h-\frac{25}{2}\Delta^\circ.
}
\]

\subsection*{Canal \(13\)}
El canal \(13\) tiene puente eléctrico directo, contraángulo de rama débil \(5\cdot4\), y torsión mínima de la diferencia octada-hexada distribuida sobre TRIT:
\[
\boxed{
\theta_{13}=A-(5\cdot4)h-\frac{8-6}{3}\Delta^\circ
=
A-20h-\frac23\Delta^\circ.
}
\]

\subsection*{Fase CP}
La fase global cierra en la rueda digital:
\[
\boxed{
\delta_{\rm CKM}=9A.
}
\]

\section{Valores}

Con
\[
A=7.297352569283801^\circ,\quad h=0.353956389828108^\circ,\quad
\Delta^\circ=0.006371085717212^\circ,
\]
se obtiene:
\[
\theta_{12}=13.003313589509^\circ,
\]
\[
\theta_{23}=2.398056157332^\circ,
\]
\[
\theta_{13}=0.213977382244^\circ,
\]
\[
\delta=65.676173123554^\circ.
\]
La matriz resultante es:
\[
|V_{\rm CKM}^{\rm HMT}|=
\begin{pmatrix}
0.974350259 & 0.225005836 & 0.003734601\\
0.224873110 & 0.973489279 & 0.041841465\\
0.008582400 & 0.041121745 & 0.999117283
\end{pmatrix}.
\]

\section{Lectura}

La corrección clave es:
\[
\boxed{
\text{los coeficientes CKM no son primero rutas de masas;}
}
\]
\[
\boxed{
\text{son primero datos de cruce }\pi\text{-Witt: }4,\ 7,\ 5,\ 6\to8,\ 3.
}
\]
El species-pack puede validar o realizar estos coeficientes en la capa física, pero la aritmética de los coeficientes ya está en el cruce:
\[
28=4\cdot7,\qquad
\frac{25}{2}=\frac{5^2}{2},\qquad
\frac23=\frac{8-6}{3},\qquad
20=5\cdot4.
\]

\section{Disciplina}

N42 es una reconstrucción formal desde los axiomas de cruce \(\pi\)-Witt. No afirma haber recuperado todavía la cita documental original perdida. Cuando reaparezcan los archivos de ruta, deberá verificarse si esta reconstrucción coincide con el protocolo histórico.
\end{document}
